\scanpage{9}\mspage{8}
$\inf$ dans $\mathcal{F}$ est $\varnothing_{\mathcal{F}}$, le plus petit élément de $\mathcal{F}$ ou \uline{figure vide}. Ainsi
\[
F\text{ et }\unsure{G|C}\text{ disjoints}\quad\Longrightarrow\quad F\text{ et }G\text{ compatibles}.
\]

On a, par les axiomes:
\[
\begin{aligned}
&&&\deleted{\text{tout}}\quad\inserted{\operatorname{omb}(F)\text{ et }\operatorname{omb}(G)}\\
 F\text{ et }G\text{ disjoints}&\quad\Longleftrightarrow\quad&&\deleted{\operatorname{omb}(F)\cap\operatorname{omb}(G)}\\
&&&\text{sont ``disjoints'' au sens}\\
&&&\text{suivant: ils sont \unsure{disjoints|disjoint}}\\
&&&\text{ensemblistement \unsure{i.e.|ie.}}\\
&&&\operatorname{omb}(F)\cap\unsure{\operatorname{om}|\operatorname{omb}}(G)=\varnothing,\ \text{et}\\
&&&\text{tout \unsure{élément|élém\textsuperscript{t}} de }\operatorname{omb}(F)\text{ est}\\
&&&\text{compatible avec tout \unsure{élt|élément}}\\
&&&\text{de }\operatorname{om}(G)
\end{aligned}\tag*{(1.13)}
\]
-- i.e. deux figures sont disjointes ssi elles n'ont pas de \unsure{lieu|lieux} communs, et \unsure{tout|$\forall$} lieu de l'une est \deleted{disjoint} \deleted{\unsure{com}} \inserted{disjoint} (au sens de \deleted{la} l'\deleted{théorie}atelier) de tout lieu de l'autre;

Cela montre que les relations de \deleted{\unsure{disj}}compatibilité et de disjonction sont importantes aussi dans les \deleted{\unsure{en de|en des}} \inserted{sous-}ensembles $\mathcal{M},\mathcal{L}$ de $\mathcal{F}$. Dans le cas des ateli\deleted{\unsure{er|ers}}ers ensemblistes, ces relations sont ``triviales'' sur $\mathcal{L}$: deux lieux sont tjrs \deleted{disjoints} \inserted{compatibles}, et disjoints ssi ils sont distincts. Par contre, même dans le cas ensembliste et $\unsure{\overline m|=}$ modéré, la relation
